\section{Methods}
\label{sec:methods}

\subsection{Cohort construction}
\label{sec:methods:cohorts}

\paragraph{Discovery cohort (TCGA).}
Formalin-fixed, paraffin-embedded H\&E-stained diagnostic whole-slide images were
obtained from The Cancer Genome Atlas via the Genomic Data Commons portal for 21
solid-tumor types (Table~S1). Slides were excluded if viable tissue
area fell below 1~mm$^2$, if severe processing artifacts were present (pen marks
covering $>$20\% of tissue area, or out-of-focus regions), or if essential
clinical metadata (vital status, follow-up time) were missing, leaving
\num{6745} diagnostic slides from \num{6088} patients (312 patients contributed
more than one slide). Slide-level descriptive analyses (the inter-measurement
correlation geometry, the UMAP embedding and pan-cancer clustering) used all
\num{6745} slides. To avoid pseudo-replication, every statistical association
used one representative slide per patient, chosen by a deterministic priority
ranking (quality-control pass status, then tissue-area adequacy, then
tile-extraction coverage, then tile count, then slide name). Twelve further
TCGA cancer types were excluded in their entirety because their dominant cell
morphologies fall outside the training domain of the cell detection model
(Table~S2). GDC file UUIDs for all \num{6745} slides are listed in
the code repository.

Matched clinical data (overall, disease-specific, disease-free and
progression-free survival; age at diagnosis, sex, pathologic stage, tissue source
site and tumor purity) came from the TCGA Pan-Cancer Clinical Data
Resource~\cite{liu2018tcgacdr}. Molecular data comprised RNA-seq expression
(RSEM-normalized), somatic mutations from the MC3 multi-caller
ensemble~\cite{ellrott2018mc3}, copy-number variation, immune cell fractions from
CIBERSORT~\cite{newman2015cibersort} and xCell~\cite{aran2017xcell}, tumor purity
from ABSOLUTE~\cite{carter2012absolute}, and immune subtypes (C1--C6) from
Thorsson et~al.~\cite{thorsson2018immune}. Molecular data were matched to slides
by TCGA case barcode (first 12 characters); for patients with multiple aliquots
the primary-tumor aliquot was used. Of \num{6745} slides, \num{5590} (82.9\%) had
matched immune subtype labels.

\paragraph{Validation cohort (CPTAC).}
Diagnostic H\&E slides and matched transcriptomic, proteomic and clinical data
were obtained for five CPTAC cancer types with TCGA counterparts: breast
carcinoma (198 slides), lung adenocarcinoma (240), head-and-neck squamous
carcinoma (390 slides from 112 patients), pancreatic adenocarcinoma (175) and
endometrial carcinoma (92), totalling \num{1095} slides from 817 patients. The
identical pipeline, model weights and feature definitions used for TCGA were
applied without retraining, recalibration or cohort-specific tuning. For patients
with multiple slides, one representative slide was selected by a deterministic
priority ranking (quality-control pass status, then tissue-area adequacy, then
tile-extraction coverage, then tile count, then fewest missing measurements, then
slide name), and all downstream analyses were performed at patient level.
Global tandem-mass-tag proteomics were available for four of the five types
(head-and-neck squamous carcinoma is a GDC-only study without a proteomic
profile).

\subsection{Tissue segmentation}
\label{sec:methods:tissue}

Tissue segmentation used a CellViT-inspired architecture~\cite{hoerst2024cellvit}
with a Phikon self-supervised ViT-B backbone~\cite{filiot2023phikon}, trained on
the PanopTILs crowdsourced annotation dataset~\cite{liu2024panoptils}. Inference was performed at
0.5~\textmu m/px on $224 \times 224$~pixel tiles with 32-pixel overlap, with the
final mask obtained by majority voting in overlap regions. The model assigns nine
tissue classes (cancerous epithelium, stroma, necrosis, normal epithelium,
tumor-infiltrating lymphocytes, junk/debris, blood, other, whitespace); mean
intersection-over-union was 0.72 on the PanopTILs held-out test set, with
substantial per-class variation and near-zero recall for necrosis and normal
epithelium in deployment, reflecting their rarity in both the training set and
resected TCGA specimens. Following the International Immuno-Oncology Biomarker
Working Group recommendation~\cite{salgado2015til}, regions classified as
tumor-infiltrating lymphocytes were reassigned to stroma before all downstream
computation, yielding five effective compartments: cancerous epithelium, stroma,
necrosis, normal epithelium and blood.

\subsection{Cell segmentation and classification}
\label{sec:methods:cells}

Cell instance segmentation and classification used
HistoPLUS~\cite{adjadj2026histoplus}, which classifies cells into 13 types; the
nine used here are
tumor cells, lymphocytes, fibroblasts, plasmocytes, neutrophils, eosinophils,
red blood cells, apoptotic bodies and mitotic figures (mean panoptic quality
0.509; per-class 0.28 to 0.73, lowest for rare types). Inference used
$224 \times 224$~pixel tiles at 40$\times$ magnification (0.25~\textmu m/px) when
available, falling back to 20$\times$ (0.50~\textmu m/px); most slides were
scanned at 40$\times$. Tiles were extracted with a 64-pixel (16~\textmu m)
overlap margin and cells in overlap regions were deduplicated by a union-find
procedure merging instances whose centroids fall within 10~\textmu m, roughly the
diameter of an epithelial nucleus. The cell model was trained on pathologist
annotations from eight cancer types (LUAD, LUSC, BRCA, COAD, BLCA, OV, PAAD,
MESO); the remaining 13 included types were processed out of distribution.
Tissue--cell discordance (cells inside a tissue-model tumor region not classified
as tumor cells) ranged from 0\% to 12\% across cancer types, highest in
pancreatic adenocarcinoma (12.2\%) despite being in distribution, plausibly
because of dense desmoplastic stroma and small tumor glands; out-of-distribution
types additionally showed inflated kinetic indices in prostate, liver and thyroid
carcinoma attributable to reduced detection sensitivity. Red blood cells and
plasmocytes were not used in measurement computation.

\subsection{Spatial compartment definition}
\label{sec:methods:spatial}

All spatial computations were performed on compartment masks resampled to a
common 8~\textmu m/px grid by nearest-neighbor interpolation, so that boundary
quantities are scanner-invariant. Connected components below
$A_{\min} = \num{2048}$~\textmu m$^2$ were removed per compartment before
distance-transform computation, to prevent segmentation fragments from inflating
boundary lengths. Five spatial bands were defined from the signed Euclidean
distance transform $d_T$ to the tumor boundary (the outer contour of the
cancerous epithelium mask; positive inside tumor): tumor front
($0 \leq d_T \leq 50$~\textmu m), tumor core ($d_T > 50$~\textmu m), near stroma
($-50 \leq d_T < 0$), far stroma ($-200 \leq d_T < -50$) and necrosis ring
(within 100~\textmu m of the necrosis boundary). The 50~\textmu m front band
corresponds to approximately five cell diameters; the 200~\textmu m stromal
cutoff is a design parameter of the pipeline. Slides were assigned a growth-pattern
regime from the tumor front fraction
$\varphi = A(B_T^{0\text{-}50}) / A(\Omega_T)$: mass-forming
($\varphi \leq 0.5$), intermediate ($0.5 < \varphi \leq 0.8$) or infiltrative
($\varphi > 0.8$). Distance-based measurements are quantized at 8~\textmu m
owing to the grid spacing.

\subsection{Histomic measurement definitions}
\label{sec:methods:features}

We computed 38 quantitative histomic measurements per slide. For presentation
they are grouped into six categories: tissue composition (3), cell densities (8),
nuclear morphology and kinetics (7), spatial organization (11), spatial
heterogeneity (5) and cell-type ratios (4). The atlas's machine-readable feature
metadata uses a coarser five-way grouping of the same 38 measurements
(composition 3, density 6, morphology 8, spatial 18, heterogeneity 3); the two
schemes partition the identical set and no measurement differs between them.
Complete definitions, units and both category assignments are given in
Table~S3.

Two candidate measurements (normal epithelium area fraction and tumor--normal
contact fraction) were removed from the atlas entirely because the tissue model
did not reliably detect normal epithelium and both showed zero signal across all
slides. Seven ratio measurements were susceptible to unstable values when
denominators approach zero; these were controlled by winsorization
(\S\ref{sec:methods:preprocessing}).

\subsection{Measurement preprocessing}
\label{sec:methods:preprocessing}

Preprocessing proceeded in three steps. First, 22 measurements with heavy
right-skew (all densities, ratios, distances and heterogeneity measures) were
log-transformed as $\log(1 + x)$. Second, all 38 were winsorized at the 0.5th and
99.5th percentiles computed per measurement across the full cohort; global rather
than per-cancer percentiles were used because the target artifacts (segmentation
failures, ratio instabilities) occur across cancer types and because several
cohorts are too small (for example cholangiocarcinoma, $n = 38$) for reliable
percentile estimation. Third, measurements were $z$-score standardized. The scope
of standardization varied by analysis: UMAP and $K$-means used global $z$-scores;
Cox regression used per-cancer-type $z$-scores, so hazard ratios are per standard
deviation within cancer type; Spearman correlations used no standardization,
being rank-based. Missing values ($<$2\% of entries) were imputed with the
column median before standardization. Raw values are retained for display in the
web interface; all models consumed preprocessed values.

\subsection{Survival analysis}
\label{sec:methods:survival}

Associations between each measurement and four endpoints (overall survival,
designated primary; disease-specific, disease-free and progression-free survival)
were assessed by Cox proportional hazards regression using
\texttt{lifelines}~\cite{davidsonpilon2019} without penalization. Each model
tests a single measurement plus covariates; no multivariate models combining
several measurements are included in the atlas. Two specifications were fitted
per measurement, cancer type and endpoint: unadjusted (measurement only) and
adjusted (measurement, age at diagnosis, sex and pathologic stage, stratified by
tissue source site). Stratification lets each site carry its own baseline hazard
without consuming degrees of freedom, the standard approach for multi-center
data~\cite{therneau2000modeling}. Pan-cancer models additionally stratified on
cancer type, estimating one shared coefficient across per-type baseline hazards.
Hazard ratios are reported per one standard deviation of the preprocessed
measurement, with 95\% Wald confidence intervals from the observed Fisher
information.

For adjusted models, missing covariates were handled by multiple imputation by
chained equations (\texttt{scikit-learn} \texttt{IterativeImputer} with Bayesian
ridge, 5 imputations, maximum 10 iterations, seed incremented per imputation),
pooled by Rubin's rules, with the fraction of missing information recorded. When
covariate missingness exceeded 20\% in any variable, total sample size fell below
30, or fewer than 3 of 5 imputations converged, complete-case analysis was used
instead. Missingness arose mainly from incomplete staging (Table~S6). A minimum of $n \geq 30$ with $\geq$10 events was required, and models
were fitted only when residual degrees of freedom remained positive.

The proportional-hazards assumption was tested for every fitted model using
Schoenfeld residuals with the Kaplan--Meier time transform~\cite{grambsch1994},
flagged as pass ($P \geq 0.05$), warn ($0.01 \leq P < 0.05$) or fail
($P < 0.01$). When the assumption was violated ($P < 0.05$, covering both warn
and fail), the hazard ratio, confidence interval and $P$ value were invalidated,
and the $P$ value was set to 1.0 for correction so that family size is preserved
without inflating discoveries. Of \num{5623} fitted models, \num{5283} yielded
valid estimates. Restricted mean survival time (RMST) differences were computed
as an assumption-free complement, as the area under the Kaplan--Meier curve to a
cancer-type-specific horizon (default \num{1095}~days; 2~years for pancreatic
adenocarcinoma and mesothelioma; 5~years for thyroid and prostate carcinoma;
Table~S7) with Irwin variance~\cite{irwin1949}, tested by
\num{5000}-permutation inference on median-split groups with \num{1000}-resample
percentile bootstrap intervals, requiring $n \geq 20$ with $\geq$5 per group. We
note that dichotomization for RMST discards information relative to continuous
methods~\cite{royston2006}. Cluster-level survival used two-sided log-rank tests
of members against non-members.

\subsection{Molecular correlations}
\label{sec:methods:correlations}

Spearman rank correlations were computed between each of the 38 measurements and
322 molecular targets: 133 curated genes drawn from OncoKB~\cite{oncokb} and the
COSMIC Cancer Gene Census~\cite{sondka2018cgc}, immune checkpoint targets and
EMT/stemness markers~\cite{nieto2016emt} (Table~S4), each assessed
for expression and copy number ($133 \times 2$); the 50 MSigDB Hallmark pathway
scores~\cite{liberzon2015}; and 6
CIBERSORT immune cell fraction scores. Pathway activity was computed by
single-sample gene set enrichment analysis (\texttt{gseapy}
1.1.12)~\cite{barbie2009ssgsea, fang2023gseapy} on the batch-corrected TCGA Pan-Cancer expression
matrix from UCSC Xena~\cite{goldman2020xena}, excluding gene sets with fewer than
10 genes present after intersection.

Unadjusted models used \texttt{scipy.stats.spearmanr} with analytical
$t$-approximation $P$ values, chosen over permutation to avoid $P$-value banding,
with 95\% percentile bootstrap intervals from \num{1000} resamples.
Covariate-adjusted models used partial Spearman correlations computed by
rank-transforming measurement, target and covariates, residualizing the ranked
measurement and target against ranked covariates by ordinary least squares, and
correlating the residuals~\cite{conover1999}; $P$ values used the standard
partial-correlation $t$ statistic with $\mathrm{df} = n - 2 - k$ for $k$
one-hot-encoded covariate columns. Tissue source site entered as a grouped
covariate (five most frequent sites per cancer type, remainder collapsed to
``Other'') rather than a stratification variable, since rank residualization
requires explicit columns. A cancer type was analyzed if $n \geq 30$; each pair
required $n \geq 10$ non-missing observations. Per-task random seeds were derived
deterministically from a hash of the cancer type and feature names. The analysis
comprised \num{487638} evaluable pairs (measurements $\times$ 322 targets
$\times$ 22 cohorts $\times$ 2 models, excluding combinations with insufficient
data). Of these, \num{461130} involve the 38 measurements; the remaining
\num{26508} involve three slide-level quality-control variables (missingness
flag, number of missing measurements and number of outlying measurements) that
were computed alongside the measurements, are released in the same table, and
entered the same correction families.

\subsection{Categorical associations}
\label{sec:methods:categorical}

Associations with categorical molecular variables (somatic mutation status,
copy-number alteration, immune subtype) were tested per cancer type. Unadjusted
two-group comparisons used the two-sided Mann--Whitney $U$ test with Cliff's
$\delta$ as effect size, computed as the mean of $\mathrm{sign}(x_i - y_j)$ over
all pairs with \num{1000}-resample bootstrap intervals (bootstrap rather than
analytical intervals because ties violate the continuity assumption of the
analytical variance formula~\cite{cliff1993}). Variables with more than two
levels used the Kruskal--Wallis $H$ test with
$\eta^2_H = (H - k + 1)/(N - k)$~\cite{tomczak2014}, floored at zero.
Covariate-adjusted models used rank-ANCOVA with Freedman--Lane permutation
inference~\cite{freedmanlane1983} (\num{1000} permutations of covariate-only
residuals), with partial $\eta^2$ as effect size. Group ordering followed a fixed
convention (mutant before wild-type, altered before neutral) so that the sign of
Cliff's $\delta$ is interpretable across analyses. A minimum of $n \geq 30$ total
and $\geq$5 per group was required.

\subsection{Clustering and embedding}
\label{sec:methods:clustering}

Unsupervised clustering was performed at two levels. At the pan-cancer (L1)
level, all \num{6745} slides were clustered on the full 38-dimensional
preprocessed vector; at the cancer-specific (L2) level, clustering was performed
within each cancer type using cancer-type-specific $z$-scores. We used $K$-means
(\texttt{scikit-learn}, $n_{\mathrm{init}} = 10$, random state 42). For L1,
solutions were computed for $K \in \{3, \ldots, 25\}$ and scored by silhouette,
Calinski--Harabasz, Davies--Bouldin and gap statistics; $K = 10$ was selected on
convergence of silhouette (local maximum), Calinski--Harabasz (plateau) and
Davies--Bouldin (local minimum), the gap statistic favouring $K = 8$. No
dimensionality reduction was applied before clustering, to keep cluster feature
profiles interpretable; the effective dimensionality of 18 components at 90\%
explained variance implies moderate redundancy, which gives correlated
measurements proportionally more weight. For L2, $K$ was selected by the elbow
method over $K \in [2, 8]$~\cite{thorndike1953}, yielding 69 subclusters; cancer
types with fewer than 20 samples received a single cluster and individual
clusters with fewer than 10 samples were dropped. Stability was assessed by 50
iterations of 80\% subsampling without replacement, re-clustering and comparing
to the original labels by adjusted Rand index and mean best-match Jaccard
index~\cite{hennig2007}. Cluster names were generated from a fixed schema
combining an immune axis label, a stromal axis label, the most extreme remaining
measurement, and a cancer-enrichment tag where one type constitutes $\geq$40\% of
the cluster. Two-dimensional visualization used UMAP
($n_{\mathrm{neighbors}} = 15$, $\mathrm{min\_dist} = 0.1$, Euclidean, random
state 42)~\cite{mcinnes2018} for display only; no inference was drawn from
embedding coordinates and all clustering was performed in the original
38-dimensional space.

\subsection{Continuous histomic programs}
\label{sec:methods:programs}

Sparse principal component analysis (\texttt{scikit-learn} \texttt{SparsePCA},
6 components, $\alpha = 1.0$, random state 42) was fitted on standardized
measurements from \num{5371} TCGA patients with complete feature vectors, after
representative-slide selection. Programs were named automatically from their
top-three loading measurements and interpreted post hoc by Spearman correlation
against all 50 Hallmark pathway scores in TCGA. The fitted scaler and
decomposition were then applied unchanged to 297 CPTAC patients with complete
feature vectors. Distributional transfer was assessed by two-sample
Kolmogorov--Smirnov tests between TCGA and CPTAC program scores. Program--protein
associations were computed as Spearman correlations against CPTAC tandem-mass-tag
protein abundance (\num{2220} program--protein pairs across four cancer types).

\subsection{External validation in CPTAC}
\label{sec:methods:validation}

\paragraph{Measurement concordance.}
For each matched cancer-type pair we computed the standardized difference
$d = (\mathrm{median}_{\mathrm{CPTAC}} - \mathrm{median}_{\mathrm{TCGA}}) /
\mathrm{SD}_{\mathrm{TCGA}}$ per measurement over all slides of that type in each
cohort (Figure~\ref{fig:fig4}A shows the signed value). A measurement passes at
$|d| \leq 1.0$; a cancer type passes when more than 80\% of measurements
($\geq 31/38$) pass. The inter-measurement correlation geometry was compared by
computing the full $38 \times 38$ Spearman matrix separately in each cohort and
correlating the flattened upper triangles (703 pairs) by Pearson correlation,
with a prespecified pass criterion of $r > 0.80$.

\paragraph{Association replication.}
Ten histomic--molecular associations were prespecified from the TCGA analysis
before CPTAC results were examined. Their CPTAC $P$ values were corrected within
the full CPTAC correction family (approximately \num{24700} expression tests per
model per cancer type), not within the prespecified subset, so replication is
assessed against the same multiplicity burden as discovery. Direction
concordance and significance replication are reported with 95\% Clopper--Pearson
exact binomial intervals. Aggregate transfer was quantified by extracting the top
50 correlated genes per measurement per cohort for 8 representative measurements
across 4 matched cancer types, matching by gene symbol, and correlating the two
cohorts' Spearman coefficients across the resulting 750 triples by Pearson
correlation with Fisher $z$-transform intervals.

\paragraph{Survival replication.}
Eight prespecified survival hypotheses were tested. Following ICH E9 guidance for
confirmatory tests of prespecified hypotheses, nominal (uncorrected) $P$ values
are reported and no false-discovery correction was applied. Direction concordance
is computed only over associations that were significant in TCGA and for which
CPTAC had sufficient follow-up. Survival follow-up was absent or minimal in
breast (0 patients with overall survival data), lung adenocarcinoma (25) and
endometrial carcinoma (8), leaving pancreatic adenocarcinoma ($n = 110$) and
head-and-neck squamous carcinoma ($n = 41$) as the only evaluable cohorts; their
minimum detectable effects at 80\% power are HR~$\geq 1.71$ and HR~$\geq 2.45$
respectively, so this analysis is reported as underpowered by construction.

\paragraph{Cluster-structure transfer.}
A \texttt{StandardScaler} and $K$-means model ($K = 10$) fitted on TCGA were
applied unchanged to CPTAC by nearest-centroid prediction. Per-cluster pathway
enrichment (Cliff's $\delta$ from Mann--Whitney $U$ tests, in-cluster versus
out-of-cluster) was computed independently in each cohort and compared by
Spearman correlation, both per cluster and across all 450 matched
cluster--pathway pairs. Adjusted Rand index, normalized mutual information and
V-measure between cluster labels and cancer type were computed in both cohorts.

\subsection{Proteomics and protein-specific signal}
\label{sec:methods:proteomics}

Spearman correlations were computed between each measurement and CPTAC global
tandem-mass-tag protein abundance in the four cancer types with proteomic
profiles. Direction concordance between messenger RNA and protein was assessed
over pairs significant in both modalities after correction, with
Clopper--Pearson intervals. To determine how much of the protein signal is
independent of transcript abundance, we computed partial Spearman correlations
between each measurement and each protein adjusting for the matched transcript,
using the same rank-residualization procedure as
\S\ref{sec:methods:correlations}, over \num{14687} (measurement, gene, cancer
type) triplets, with Benjamini--Hochberg correction applied per cancer type.
Retention was summarized as the fraction of triplets for which the partial
coefficient retained more than half the magnitude of the unadjusted coefficient.

\subsection{Added-value analyses}
\label{sec:methods:addedvalue}

To test whether measurements carry information beyond bulk transcriptomics, Cox
models for four target measurements were refitted with matched molecular
covariates added to the clinical specification: \textit{CD8A}, \textit{TIGIT},
immune subtype and leukocyte fraction for intratumoral lymphocyte density and
tumor--lymphocyte nearest-neighbor distance;
\textsc{hallmark\_g2m\_checkpoint} and \textit{MKI67} for mitotic index; and
\textit{ESR1} plus PAM50 subtype for nuclear pleomorphism in breast carcinoma.
Incremental discrimination was quantified as Harrell's concordance index for four
nested models per cancer type (clinical only; clinical + RNA; clinical +
histomics; clinical + RNA + histomics), summarized as the median across the 18
cancer types with evaluable models, together with the fraction of cancer types in
which the increment was positive. Concordance was evaluated on the data used for
fitting, so the increments are apparent rather than cross-validated and are an
upper bound on out-of-sample gain.

\subsection{Multiple testing correction}
\label{sec:methods:multiple}

All $P$ values were corrected by the Benjamini--Hochberg
procedure~\cite{benjamini1995} within explicitly declared families, chosen to
pool biologically coherent tests without the conservatism of pooling unrelated
analyses (Table~S5 gives every family definition). Survival
families were defined per cancer type $\times$ endpoint $\times$ model, so that
correction pools only the 38 measurement-level tests within one biological
context; correlation families were defined per cancer type $\times$ molecular
data type $\times$ model, giving 150 families. Every adjusted $P$ value is stored
alongside its correction-family identifier and family size, so that any reported
result can be re-derived. Gene set enrichment analysis used its own canonical
false-discovery procedure based on pooled null normalized enrichment score
distributions~\cite{subramanian2005} rather than Benjamini--Hochberg. Overall
survival is the primary endpoint; the other three are sensitivity analyses,
corrected independently per endpoint rather than pooled, because the endpoints
share events and therefore violate the independence assumption of joint
correction. Permutation $P$ values used the add-one correction of Phipson and
Smyth~\cite{phipson2010permutation}. Bootstrap intervals used \num{1000}
percentile resamples and were reported only when $\geq$90\% of resamples produced
valid estimates. Unless stated otherwise, ``significant'' means
$P_{\mathrm{adj}} < 0.05$.

\subsection{Power analysis and evidence grading}
\label{sec:methods:power}

The minimum detectable effect size at 80\% power ($\alpha = 0.05$, two-sided) was
computed for every association: by the Schoenfeld--Freedman
approximation~\cite{schoenfeld1983, freedman1982} for survival, the Fisher
$z$-transform for correlations, and simulation-based power curves (\num{2000}
simulated datasets with binary search) for categorical associations. Every
association carries an evidence badge: \textbf{strong} ($P_{\mathrm{adj}} <
0.01$, effect above the medium-effect threshold, narrow interval, $n \geq 100$),
\textbf{moderate} ($P_{\mathrm{adj}} < 0.05$, effect above the small-effect
threshold, narrow or moderate interval, $n \geq 50$), \textbf{suggestive}
($P_{\mathrm{adj}} < 0.10$ or interval excluding the null, $n \geq 30$) or
\textbf{insufficient} ($n < 30$ or missing statistics). Thresholds for $\rho$ and $\eta^2$ follow
Cohen~\cite{cohen1988}, those for Cliff's $\delta$ approximate Romano
et~al.~\cite{romano2006}, and those for hazard ratios were set for this study:
medium, HR~$\geq 1.5$ or
$\leq 0.667$, $|\rho| \geq 0.3$, $|\delta| \geq 0.3$, $\eta^2 \geq 0.06$; small,
HR~$\geq 1.18$ or $\leq 0.847$, $|\rho| \geq 0.1$, $|\delta| \geq 0.15$,
$\eta^2 \geq 0.01$.

\subsection{Quality control, batch effects and sensitivity analyses}
\label{sec:methods:qc}

Batch effects from tissue source site were assessed by principal variance
component analysis~\cite{li2009pvca}, decomposing variance in the top principal
components (explaining $\geq$80\% cumulative variance, up to 10 components) into
site, cancer type and residual fractions using marginal one-way ANOVA $\eta^2$
weighted by explained variance. Because site and cancer type are not orthogonal,
marginal $\eta^2$ values can exceed 1.0 in sum; residual variance was computed as
$\max(0, 1 - \eta^2_{\mathrm{batch}} - \eta^2_{\mathrm{bio}})$ and all three
fractions renormalized. Silhouette analysis treated site labels as cluster
assignments on standardized measurements (Euclidean, subsampled to \num{5000}
slides), with scores above 0.25 flagged as moderate and above 0.5 as strong batch
effects.

Four sensitivity analyses were run. (i) A permutation null shuffled molecular
labels within cancer type and re-ran the full production correction pipeline over
the \num{39900} measurement $\times$ pathway $\times$ cancer-type tests for 100
permutations. (ii) The Benjamini--Yekutieli procedure, valid under arbitrary
dependence, was substituted for Benjamini--Hochberg; for this comparison all
tests of a given analysis type were pooled into a single family, so its
denominator (272 significant survival associations) differs from the production
149-family design (260). (iii) Top correlations were recomputed with molecular
subtype as an additional covariate (PAM50 for breast, molecular class for
endometrial, immune subtype for pancreatic carcinoma; the lung adenocarcinoma
driver stratification was uninformative because neither \textit{EGFR} nor
\textit{KRAS} mutation was prevalent enough to form contrasting groups).
(iv) Survival models for the top 10 measurements were refitted with leukocyte
fraction as an additional covariate.

\subsection{Software, reproducibility and interfaces}
\label{sec:methods:software}

All analyses were implemented in Python 3.11 and orchestrated as a directed
acyclic graph by Snakemake~\cite{moeldner2021snakemake}, which supports a
subsetted dry-run mode for validating the full pipeline before production runs.
Key versions: \texttt{lifelines}~0.29.0, \texttt{scipy}~1.12.0,
\texttt{scikit-learn}~1.4.0, \texttt{umap-learn}~0.5.5, \texttt{gseapy}~1.1.12,
\texttt{statsmodels}~0.14.1, \texttt{numpy}~1.26.4, \texttt{pandas}~2.2.0. Every
stochastic procedure (permutation, bootstrap, subsampling, $K$-means
initialization, UMAP) used an explicit seed; per-task seeds in embarrassingly
parallel stages were derived deterministically from task identifiers so that
results do not depend on execution order. Each run writes a provenance record
containing pipeline version, timestamps, input checksums, an output manifest and
an automated validation report.

The web interface is an Astro static site with React components; all statistical
results, feature profiles, cluster metadata and visualization data are serialized
to static JSON at build time, so no analysis server runs at request time. Three
interpretability layers connect statistics to tissue: a per-slide compartment map
rendering the five tissue compartments across nine spatial zones; for each
measurement, the top-ranked $224 \times 224$~pixel tiles scored by that
measurement's own computation strategy, each with slide coordinates, a scale bar
and a toggleable cell-type overlay. The overlay renders the full set of classes
the cell model emits (14 in the current release) for visual inspection, whereas
the nine classes listed in \S\ref{sec:methods:cells} are those used to compute
the measurements; the remaining classes are displayed but not quantified. The
interface also provides bidirectional links from association
tables to measurement pages and from measurement pages to specific tiles on
specific slides. The same tables are exposed programmatically through 45 typed
tools over a Model Context Protocol server, so that automated agents can issue
the same queries as human users.

\subsection{Figure preparation}
\label{sec:methods:figures}

Every figure panel was drawn from data exported directly from the frozen pipeline
outputs (TCGA run of 9 March 2026; CPTAC run of 17 March 2026) by a single script
(\texttt{patterns\_research\_article/scripts/export\_figure\_data.py}), which writes
one machine-readable file per figure containing every plotted value and every
statistic quoted in the legends. A second script
(\texttt{build\_figures.py}) renders these files deterministically as vector
graphics at print size, and a third (\texttt{export\_figures.py}) produces the
submitted PDF and TIFF files; rerunning the three scripts regenerates every
figure from the released tables. Hematoxylin-and-eosin tiles and cell-call overlays in
Figures~\ref{fig:fig1}A and \ref{fig:fig6} are the pipeline's own image files,
reproduced without alteration other than uniform scaling; compartment maps are
label images in which the pipeline's class colors were remapped to the figure
palette. Dense scatter plots are shown as two-dimensional histograms with
logarithmic color scaling, and point clouds in Figure~\ref{fig:fig3}C are shown on
percentile-rank scales, on which a Spearman coefficient is the Pearson coefficient
of the plotted points.
